\chapter{Maps, Objectives and the Two Authors}

\section{The Anatomy of a PGM Map}

PGM is a server plugin that runs competitive Minecraft matches. It loads a map,
which is two things bundled together: a Minecraft world, meaning the blocks, the
terrain and the buildings; and a file called \id{map.xml} that says what the world
is \emph{for}. Without the XML the world is scenery. The XML is what states that
this team spawns here, that this region cannot be built in, that the match ends
when that block is broken.

Objectives come in a handful of kinds, and two of them carry this paper. A
\textbf{Capture the Wool} map, universally abbreviated CTW, hides coloured wool in
rooms, and a team wins by carrying the right wool back to its own monument. A
\textbf{destroy-the-monument} map, DTM, puts a block or a small structure in each
team's ground, and a team wins by breaking the enemy's. Others exist: a core that
leaks lava when broken, a control point held over time. The studio builds them
all, but CTW and DTM are the two worked examples here.

\section{Team Counts by Objective}
\label{sec:team-counts}

The number of teams a map can hold is not one question but two, and they have
different answers. One is about the objective and belongs to PGM. The other is
about symmetry and belongs to the studio.

Capture the Wool works at two teams or four, and four-team CTW maps have been
built. It works at four because each team has a \emph{distinct} objective:
capturing the other three teams' wools. Control points, whose support is recent,
take four teams or more, usually an even number, simply because maps are usually
symmetric.

Destroyables and cores are different, and the reason is worth stating properly
because it is easy to mistake for a software limit. It is not one. The difficulty
is \textbf{ownership}. A destroyable cannot be assigned to a team the way a wool
can. If a team's monument is destroyed then that team has, in a sense, lost. But
that does not determine a \emph{winner}, and with more than two teams there is no
well-defined one. PGM maps work around it in two ways. Some make the match blitz,
so the last surviving team wins. Others move the players of a team that has lost
its monument onto another team, so that one team can truly win. Neither of those
is raw DTM or DTC. A destroy map without one of those modes layered on it is
therefore a two-team map in practice.

\begin{ruling}{The Studio's Team Ceiling}
The studio's automatic symmetry fan tops out at four teams, because \id{rot\_90}
is the only four-image mode it has. That is what caps the studio, rather than
anything about the objectives. An agent driving the API directly can author
anything, a six-team map included. What it gives up is the automatic fan, so each
team's ground has to be authored rather than mirrored into place.
\end{ruling}

The blocks and the XML have to agree with each other, and that is harder than it
sounds. A wool room's position in the XML is a coordinate; whether that coordinate
is inside a room a player can reach is a fact about the blocks. A build region
drawn in the XML means nothing if the ground it covers was never placed. Most of
what a map-authoring program does is keeping those two descriptions true of one
another.

\section{The Studio: Scope and Interface}

\textsf{pgm-studio} is an ASP.NET Core application with a browser client, a
database, and an HTTP API of \RoutePathCount{} paths carrying \RouteCount{}
operations. It authors maps: it plans a layout, draws its ground, states what the
map is played for, and writes out the world and the \id{map.xml} a server can
load.

The browser client exists and is used. But the API is the whole surface, and every
map described in this paper was made without opening the client once. A caller
with a terminal can compose a board, reshape its ground, state its objectives,
build the world and read the result back, entirely over HTTP. That property is
what makes the rest of this paper possible, and it was designed in rather than
discovered. The documentation for each tool carries that tool's endpoints, on the
principle that a description which cannot be acted on is a description of nothing.

\section{The Authoring Agent}

For most of the maps in the repository this paper draws on, the operator is a
language model running in a terminal. It reads documentation, writes a Python
script that emits two JSON documents, posts them, reads the studio's answer and
corrects. It has no eyes on the world. It cannot walk around the map. Every belief
it holds about what it has built comes from a number the studio handed back, and a
good deal of the discipline described in Part~III exists because that is a
genuinely difficult position to author from.

This paper calls that operator \agent{the agent} throughout, and the program it
drives \studio{the studio}. The two words are used precisely and always mean those
two things.

\section{The Central Argument}

It is tempting to describe a system like this as a generator with a human in the
loop, or as an editor with some automation in it. Neither is right, and the reason
neither is right is the substance of this paper.

The studio is not a generator, because it does not produce maps. Its composer
produces an \emph{arrangement}: where the hub sits, how far the wool rooms stand
off it, which side the front line faces. An arrangement is not a map. It has no
elevation, no curve, no building, no second storey, no material. Every composed
board is flat, rectangular and grey. What it buys is the one thing that is hard to
invent and easy to get wrong, and the repository's own brief on the subject puts it
as sharply as it can be put: a composed board is a suggestion, not a plan.

The studio is not an editor either, because a very large amount of what ends up in
the world was never stated by the agent. The agent says a theme should follow the
ground's slope, and the studio decides which of five materials lands on each of
several hundred thousand cells. The agent says a house stands here, and the studio
decides its storey heights, its roof pitch, where its door goes and which way the
door faces. The agent says the two halves of the board are mirror images, and the
studio decides what that means for every shape, every prop and every region. These
are not defaults in the sense of values that could have been typed out. They are
derivations, and the agent could not have produced them by hand in any reasonable
time.

The honest description is therefore the third one: two parties, a boundary between
them, and the boundary in a different place at every stage. At the board level the
studio arranges and the agent adapts. At the ground level the agent states intent
and the studio computes geometry. At the paint level the agent picks a rule and the
studio applies it. At the objective level the agent decides almost everything, and
the studio's task is to refuse the decisions that cannot be built.

That last one deserves its own sentence, because it is the part least visible from
outside. A large fraction of the studio consists of gates: \RuleCount{} rules
across \RuleFamilyCount{} families, each of which can stop a request and say why.
They are not validation in the defensive sense. They are the accumulated knowledge
of what makes a map fail, expressed as things the software will not let happen, and
an agent authoring against them is being taught as much as it is being checked.

\section{Structure of This Paper}

The paper works through the stack once, in the order a map passes through it, and
at each stage says which party decided what. Then it builds two maps.

The first is a CTW board in the \emph{micro} size band, pulled straight off the
composer and then adapted: reshaped, given elevation, given buildings, given a
second storey, until it is a place rather than a diagram. The second is a DTM
board authored from nothing, with no composed arrangement to start from, by an
agent given no brief beyond the objective. Both are shown at every stage of the
pipeline, from the bare rectangles the composer emitted to the finished world.

The last chapter is about the agent rather than the map. While the DTM board was
being authored, every HTTP call its author made was recorded: the route, the
status, the rules that fired, the time it took. That recording is what the chapter
is written from, so the account of how an authoring session goes is taken off the
wire rather than from the agent's own description of itself.
